\begin{table}[H]
	%\begin{table}[ht]
	\caption{History}
	\label{tab:arch_history03}
	\centering
	\begin{tabularx}{\textwidth}{|X |X |X |}
	  \hline
	  \multicolumn{3}{|l|}{History} \\
	  \hline
	  \multicolumn{1}{|l}{Target Spec.} & \multicolumn{1}{l}{Current version : 1.0, 2025-02-22} & \multicolumn{1}{l|}{Author: A. Siggelkow}\\
	  \multicolumn{1}{|l}{ } & \multicolumn{1}{l}{Previous version : 0.0, 2020-10-12} & \multicolumn{1}{l|}{ }\\
	  \hline
	  \hline
	  Paragraph (in previous version) & Paragraph (in current version) / date & Subjects (major changes since last revision) \\
	  \hline
	  \hline
	  \multicolumn{3}{|l|}{Current Chapter Version: C0.1; Date: 2025-02-22} \\
	  \hline
	  & 2025-02-22 & First draft (AS) \\
	  \hline
	\end{tabularx}
\end{table}

\subsection{Required Control and Status Registers (CSRs)}
	
	\subsubsection{Machine Trap Setup CSRs}
		\begin{itemize}
			\item \textbf{mstatus (0x300):} Machine status register.
			\item \textbf{misa (0x301):} ISA and extensions
			\item \textbf{medeleg (0x302):} Machine exception delegation register.
			\item \textbf{mideleg (0x303):} Machine interrupt delegation register.
			\item \textbf{mie (0x304):} Machine interrupt-enable register.
			\item \textbf{mtvec (0x305):} Machine trap-handler base address.
			\item \textbf{mcounteren (0x306):} Machine counter enable.
		\end{itemize}
	
	\subsubsection{Machine Trap Handling CSRs}
		\begin{itemize}
			\item \textbf{mscratch (0x340):} Machine scratch register.
			\item \textbf{mepc (0x341):} Program counter saved at the time of the trap.
			\item \textbf{mcause (0x342):} Reason for the trap (Interrupt/Exception code).
			\item \textbf{mtval (0x343):} Machine trap value.
			\item \textbf{mip (0x344):} Machine interrupt pending.
			\item \textbf{mtinst (0x34A):} Machine trap instruction (transformed
			\item \textbf{mtval2 (0x34B):} Machine second trap value.
		\end{itemize}
	
	\subsubsection{Supervisor Trap Setup CSRs}
		\begin{itemize}
			\item \textbf{sstatus (0x100):} Restricted view/alias of \textit{mstatus}.
			\item \textbf{sie (0x104):} Machine exception delegation register.
			\item \textbf{stvec (0x105):} Base address for S-mode trap handlers.
			\item \textbf{scounteren (0x106):} Supervisor counter enable.
		\end{itemize}
	
	\subsubsection{Supervisor Trap Handling CSRs}
		\begin{itemize}
			\item \textbf{sscratch (0x140):} Supervisor scratch register.
			\item \textbf{sepc (0x141):} Saved program counter for S-mode.
			\item \textbf{scause (0x142):} Supervisor trap cause.
			\item \textbf{stval (0x143):} Supervisor trap value
			\item \textbf{sip (0x144):} Supervisor interrupt pending.
			\item \textbf{scountovf (0xDA0):} Supervisor count overflow.
		\end{itemize}

\subsection{Mascine mode CSR desriptions} % cap 3
	\subsubsection{Machine Status (mstatus and mstatush) Registers}
		The mstatus register tracks and controls the current operating state of the hart, including global interrupt enables and privilege levels.
		It is crucial for context switching because it preserves previous states during transitions to and from trap handlers.
		Additionally, it manages the status of architectural extensions like floating-point and vector units.
		
		
		\begin{figure}[!htbp]
			\centering
			\begin{bytefield}[bitwidth=0.06\textwidth, bitheight=3ex]{16}
				
				% ROW 2: Bits 63 to 48
				\bitbox[]{1}{\tiny 63} & \bitbox[]{1}{\tiny 62} & \bitbox[]{13}{} & \bitbox[]{1}{\tiny 48} \\
				\bitbox{1}{\tiny SD} & \bitbox{15}{\tiny WPRI} \\
				\\[-2ex]
				
				% ROW 2: Bits 47 to 32
				\bitbox[]{1}{\tiny 47} & \bitbox[]{3}{} & \bitbox[]{1}{\tiny 43} & \bitbox[]{1}{\tiny 42} & \bitbox[]{1}{\tiny 41} & \bitbox[]{1}{\tiny 40} & \bitbox[]{1}{\tiny 39} & \bitbox[]{1}{\tiny 38} & \bitbox[]{1}{\tiny 37} & \bitbox[]{1}{\tiny 36} & \bitbox[]{1}{\tiny 35} & \bitbox[]{1}{\tiny 34} & \bitbox[]{1}{\tiny 33} & \bitbox[]{1}{\tiny 32} \\
				\bitbox{5}{\tiny WPRI} & \bitbox{1}{\tiny MDT} & \bitbox{1}{\tiny MPELP} & \bitbox{1}{\tiny WPRI} & \bitbox{1}{\tiny MPV} & \bitbox{1}{\tiny GVA} & \bitbox{1}{\tiny MBE} & \bitbox{1}{\tiny SBE} & \bitbox{2}{\tiny SXL[1:0]} & \bitbox{2}{\tiny UXL[1:0]} \\
				\\[-2ex]
				
				% ROW 3: Bits 31 to 16
				\bitbox[]{1}{\tiny 31} & \bitbox[]{5}{} & \bitbox[]{1}{\tiny 25} & \bitbox[]{1}{\tiny 24} & \bitbox[]{1}{\tiny 23} & \bitbox[]{1}{\tiny 22} & \bitbox[]{1}{\tiny 21} & \bitbox[]{1}{\tiny 20} & \bitbox[]{1}{\tiny 19} & \bitbox[]{1}{\tiny 18} & \bitbox[]{1}{\tiny 17} & \bitbox[]{1}{\tiny 16} \\
				\bitbox{7}{\tiny WPRI} & \bitbox{1}{\tiny SDT} & \bitbox{1}{\tiny SPELP} & \bitbox{1}{\tiny TSR} & \bitbox{1}{\tiny TW} & \bitbox{1}{\tiny TVM} & \bitbox{1}{\tiny MXR} & \bitbox{1}{\tiny SUM} & \bitbox{1}{\tiny MPRV} & \bitbox{1}{\tiny XS[1:0]} \\
				\\[-2ex]
				
				
				% ROW 4: Bits 15 to 0
				\bitbox[]{1}{\tiny 15} & \bitbox[]{1}{\tiny 14} & \bitbox[]{1}{\tiny 13} & \bitbox[]{1}{\tiny 12} & \bitbox[]{1}{\tiny 11} & \bitbox[]{1}{\tiny 10} & \bitbox[]{1}{\tiny 9} & \bitbox[]{1}{\tiny 8} & \bitbox[]{1}{\tiny 7} & \bitbox[]{1}{\tiny 6} & \bitbox[]{1}{\tiny 5} & \bitbox[]{1}{\tiny 4} & \bitbox[]{1}{\tiny 3} & \bitbox[]{1}{\tiny 2} & \bitbox[]{1}{\tiny 1} & \bitbox[]{1}{\tiny 0} \\
				\bitbox{1}{\tiny XS[1:0]} & \bitbox{2}{\tiny FS[1:0]} & \bitbox{2}{\tiny MPP[1:0]} & \bitbox{2}{\tiny VS[1:0]} & \bitbox{1}{\tiny SPP} & \bitbox{1}{\tiny MPIE} & \bitbox{1}{\tiny UBE} & \bitbox{1}{\tiny SPIE} & \bitbox{1}{\tiny WPRI} & \bitbox{1}{\tiny MIE} & \bitbox{1}{\tiny WPRI} & \bitbox{1}{\tiny SIE} & \bitbox{1}{\tiny WPRI}
				
			\end{bytefield}
			\caption{Machine-mode status (\texttt{mstatus}) register for RV64.}
		\end{figure}
	
	\subsubsection{Machine Trap-Vector Base-Address (mtvec) Register}
		The mtvec register holds the base address and mode configuration used by the hardware when a trap enters Machine mode.
		It supports both Direct mode for single-address jumps and Vectored mode for specific interrupt offsets.
		Proper setup is essential for ensuring the system can accurately respond to exceptions and hardware events.
		
		\begin{figure}[!htbp]
			\centering
			\begin{bytefield}[bitwidth=0.06\textwidth, bitheight=4ex]{16}
				
				% HEADER
				\bitbox[]{2}{\raggedright \tiny MXLEN-1} & \bitbox[]{11}{} & \bitbox[]{1}{\tiny 2} & \bitbox[]{1}{\tiny 1} & \bitbox[]{1}{\tiny 0} \\
				
				% BODY
				\bitbox{14}{\tiny BASE[MXLEN-1:2] (WARL)} & \bitbox{2}{\tiny MODE (WARL)} \\
				
				% FOOTER
				\bitbox[]{14}{\tiny MXLEN-2} & \bitbox[]{2}{\tiny 2}
				
			\end{bytefield}
			\caption{Encoding of \texttt{mtvec} MODE field.}
		\end{figure}
	
	\subsubsection{Machine Trap Delegation (medeleg and mideleg) Registers}
		The medeleg register allows specific synchronous exceptions to be delegated to a lower privilege mode, such as Supervisor mode.
		By setting bits here, the hardware avoids unnecessary transitions into Machine mode for common events like page faults.
		This significantly reduces software overhead in systems running a full operating system.
		
		\begin{figure}[!htbp]
			\centering
			\begin{bytefield}[bitwidth=0.06\textwidth, bitheight=3ex]{16}
				
				\bitbox[]{1}{\raggedright \tiny 63} & \bitbox[]{14}{} & \bitbox[]{1}{\raggedleft \tiny 0} \\
				
				\bitbox{16}{\tiny Synchronous Exceptions (WARL)} \\
				
				\bitbox[]{16}{\tiny 64}
				
			\end{bytefield}
			\caption{Machine Exception Delegation (\texttt{medeleg}) register.}
			\label{fig:medeleg}
		\end{figure}
		
		The mideleg register functions similarly to its exception counterpart by delegating specific hardware interrupts to the Supervisor level.
		When a delegated interrupt occurs, the hardware automatically jumps to the S-mode trap handler instead of the M-mode one.
		This is a foundational feature for building efficient operating systems that handle their own timing and I/O.
		
		\begin{figure}[!htbp]
			\centering
			\begin{bytefield}[bitwidth=0.06\textwidth, bitheight=3ex]{16}
				
				\bitbox[]{2}{\raggedright \tiny MXLEN-1} & \bitbox[]{13}{} & \bitbox[]{1}{\raggedleft \tiny 0} \\
				
				\bitbox{16}{\tiny Interrupts (WARL)} \\
				
				\bitbox[]{16}{\tiny MXLEN}
				
			\end{bytefield}
			\caption{Machine Interrupt Delegation (\texttt{mideleg}) Register.}
			\label{fig:mideleg}
		\end{figure}
	
	\subsubsection{Machine Interrupt (mip and mie) Registers}
		The mip register contains bits that indicate which interrupts are currently pending at the Machine level, reflecting the real-time status of hardware signals.
		While some bits are read-only hardware indicators, others can be toggled by software to trigger inter-processor interrupts.
		Monitoring this register allows a kernel to prioritize and service various pending event requests effectively.
		
		\begin{figure}[!htbp]
			\centering
			\begin{bytefield}[bitwidth=0.06\textwidth, bitheight=3ex]{16}
				
				\bitbox[]{2}{\raggedright \tiny MXLEN-1} & \bitbox[]{13}{} & \bitbox[]{1}{\raggedleft \tiny 0} \\
				
				\bitbox{16}{\tiny Interrupts (WARL)} \\
				
				\bitbox[]{16}{\tiny MXLEN}
				
			\end{bytefield}
			\caption{Machine Interrupt-Pending (\texttt{mip}) register.}
			\label{fig:mip}
		\end{figure}
		
		The mie register provides individual enable bits for each interrupt source recognized by the hart at the Machine level.
		It works alongside the global enable bit in the status register to determine if a pending interrupt should actually be serviced.
		Controlling these bits allows for fine-grained management of system responsiveness and priority.
		
		\begin{figure}[!htbp]
			\centering
			\begin{bytefield}[bitwidth=0.06\textwidth, bitheight=3ex]{16}
				
				\bitbox[]{2}{\raggedright \tiny MXLEN-1} & \bitbox[]{13}{} & \bitbox[]{1}{\raggedleft \tiny 0} \\
				
				\bitbox{16}{\tiny Interrupts (WARL)} \\
				
				\bitbox[]{16}{\tiny MXLEN}
				
			\end{bytefield}
			\caption{Machine Interrupt-Enable (\texttt{mie}) register.}
			\label{fig:mie}
		\end{figure}
		\subsubsection{Machine Counter-Enable (mcounteren) Register}
		
		The mcounteren register determines whether hardware performance counters are accessible from lower privilege modes.
		Machine mode uses this register to permit or restrict access to cycle counts and instruction metrics for security and privacy reasons.
		This provides a primary defense against side-channel attacks that rely on high-precision hardware timing.
		
		\begin{figure}[!htbp]
			\centering
			\begin{bytefield}[bitwidth=0.06\textwidth, bitheight=3ex]{16}
				
				\bitbox[]{1}{\tiny 31} & \bitbox[]{1}{\tiny 30} & \bitbox[]{1}{\tiny 29} & \bitbox[]{1}{\tiny 28} & \bitbox[]{5}{} & \bitbox[]{1}{\tiny 6} & \bitbox[]{1}{\tiny 5} & \bitbox[]{1}{\tiny 4} & \bitbox[]{1}{\tiny 3} & \bitbox[]{1}{\tiny 2} & \bitbox[]{1}{\tiny 1} & \bitbox[]{1}{\tiny 0} \\
				
				\bitbox{1}{\tiny HPM31} & \bitbox{1}{\tiny HPM30} & \bitbox{1}{\tiny HPM29} & \bitbox{7}{\tiny ...} & \bitbox{1}{\tiny HPM5} & \bitbox{1}{\tiny HPM4} & \bitbox{1}{\tiny HPM3} & \bitbox{1}{\tiny IR} & \bitbox{1}{\tiny TM} & \bitbox{1}{\tiny CY} \\
				
				\bitbox[]{1}{\tiny 1} & \bitbox[]{1}{\tiny 1} & \bitbox[]{1}{\tiny 1} & \bitbox[]{7}{\tiny 23} & \bitbox[]{1}{\tiny 1} & \bitbox[]{1}{\tiny 1} & \bitbox[]{1}{\tiny 1} & \bitbox[]{1}{\tiny 1} & \bitbox[]{1}{\tiny 1} & \bitbox[]{1}{\tiny 1}
				
			\end{bytefield}
			\caption{Counter-enable (\texttt{mcounteren}) register.}
			\label{fig:mcounteren}
		\end{figure}
	
	\subsubsection{Machine Counter-Inhibit (mcountinhibit) Register}
		\begin{figure}[!htbp]
			\centering
			\begin{bytefield}[bitwidth=0.06\textwidth, bitheight=3ex]{16}
				
				\bitbox[]{1}{\tiny 31} & \bitbox[]{1}{\tiny 30} & \bitbox[]{1}{\tiny 29} & \bitbox[]{1}{\tiny 28} & \bitbox[]{5}{} & \bitbox[]{1}{\tiny 6} & \bitbox[]{1}{\tiny 5} & \bitbox[]{1}{\tiny 4} & \bitbox[]{1}{\tiny 3} & \bitbox[]{1}{\tiny 2} & \bitbox[]{1}{\tiny 1} & \bitbox[]{1}{\tiny 0} \\
				
				\bitbox{1}{\tiny HPM31} & \bitbox{1}{\tiny HPM30} & \bitbox{1}{\tiny HPM29} & \bitbox{7}{\tiny ...} & \bitbox{1}{\tiny HPM5} & \bitbox{1}{\tiny HPM4} & \bitbox{1}{\tiny HPM3} & \bitbox{1}{\tiny IR} & \bitbox{1}{\tiny 0} & \bitbox{1}{\tiny CY} \\
				
				\bitbox[]{1}{\tiny 1} & \bitbox[]{1}{\tiny 1} & \bitbox[]{1}{\tiny 1} & \bitbox[]{7}{\tiny 23} & \bitbox[]{1}{\tiny 1} & \bitbox[]{1}{\tiny 1} & \bitbox[]{1}{\tiny 1} & \bitbox[]{1}{\tiny 1} & \bitbox[]{1}{\tiny 1} & \bitbox[]{1}{\tiny 1}
				
			\end{bytefield}
			\caption{Counter-enable (\texttt{mcountinhibit}) register.}
			\label{fig:mcountinhibit}
		\end{figure}
		
	\subsubsection{Machine Scratch (mscratch) Register}
		The mscratch register typically holds a pointer to a hart-local context space used during Machine-mode trap handling.
		Software often swaps its contents with a general-purpose register at the start of a trap to provide a safe "scratch pad" for saving the current state.
		This mechanism ensures that the trap handler can operate without corrupting any data belonging to the interrupted program.
		
		\begin{figure}[!htbp]
			\centering
			\begin{bytefield}[bitwidth=0.06\textwidth, bitheight=3ex]{16}
				
				\bitbox[]{2}{\raggedright \tiny MXLEN-1} & \bitbox[]{13}{} & \bitbox[]{1}{\raggedleft \tiny 0} \\
				
				\bitbox{16}{\tiny mscratch} \\
				
				\bitbox[]{16}{\tiny MXLEN}
				
			\end{bytefield}
			\caption{Machine-mode scratch register.}
			\label{fig:mscratch}
		\end{figure}
	
	\subsubsection{Machine Exception Program Counter (mepc) Register}
		The mepc register automatically saves the address of the interrupted instruction whenever a trap is taken into Machine mode.
		It serves as the target address for the mret instruction, allowing the processor to resume execution exactly where it stopped.
		Software can also manually modify this value to redirect program flow after handling a specific exception.
		
		\begin{figure}[!htbp]
			\centering
			\begin{bytefield}[bitwidth=0.06\textwidth, bitheight=3ex]{16}
				
				\bitbox[]{2}{\raggedright \tiny MXLEN-1} & \bitbox[]{13}{} & \bitbox[]{1}{\raggedleft \tiny 0} \\
				
				\bitbox{16}{\tiny mepc} \\
				
				\bitbox[]{16}{\tiny MXLEN}
				
			\end{bytefield}
			\caption{Machine exception program counter register.}
			\label{fig:mepc}
		\end{figure}
	
	\subsubsection{Machine Cause (mcause) Register}
		The mcause register identifies the specific reason for a trap by storing an interrupt bit and a unique exception code.
		It allows the trap handler to quickly differentiate between timer interrupts, external interrupts, and various synchronous exceptions.
		This information is the primary key used by software to dispatch the correct service routine for the event.
		
		\begin{figure}[!htbp]
			\centering
			\begin{bytefield}[bitwidth=0.06\textwidth, bitheight=3ex]{16}
				
				\bitbox[]{2}{\raggedright \tiny MXLEN-1} & \bitbox[]{2}{\raggedright \tiny MXLEN-2} & \bitbox[]{11}{} & \bitbox[]{1}{\raggedleft \tiny 0} \\
				
				\bitbox{2}{\tiny Interrupt} & \bitbox{14}{\tiny Exception Code (WLRL)} \\
				
				\bitbox[]{2}{\tiny 1} & \bitbox[]{14}{\tiny MXLEN-1}
				
			\end{bytefield}
			\caption{Machine Cause (\texttt{mcause}) register.}
			\label{fig:mcause}
		\end{figure}
	
	\subsubsection{Machine Trap Value (mtval) Register}
		The mtval register provides auxiliary information about a trap, such as the faulting memory address during a page fault.
		This specific context is vital for the operating system to diagnose the failure and take appropriate corrective action.
		It complements the cause register by adding the necessary "where" or "what" to the "why" of the trap.
		
		\begin{figure}[!htbp]
			\centering
			\begin{bytefield}[bitwidth=0.06\textwidth, bitheight=3ex]{16}
				
				\bitbox[]{2}{\raggedright \tiny MXLEN-1} & \bitbox[]{13}{} & \bitbox[]{1}{\raggedleft \tiny 0} \\
				
				\bitbox{16}{\tiny mtval} \\
				
				\bitbox[]{16}{\tiny MXLEN}
				
			\end{bytefield}
			\caption{Machine Trap Value (\texttt{mtval}) register.}
			\label{fig:mtval}
		\end{figure}

\subsection{Machine mode CSR desriptions}	% cap 12
	\subsubsection{Memory Privilege in sstatus Register}
		The sstatus register is the Supervisor-level version of the status register, tracking the hart’s state during S-mode execution.
		It mirrors several fields from the Machine-mode status register but is restricted to functionality relevant to the operating system kernel.
		Kernels use it to manage global interrupts and track previous privilege states during trap returns.
		
		\begin{figure}[!htbp]
			\centering
			\begin{bytefield}[bitwidth=0.06\textwidth, bitheight=3ex]{16}
				
				\bitbox[]{1}{\tiny 63} & \bitbox[]{1}{\tiny 62} & \bitbox[]{13}{} & \bitbox[]{1}{\tiny 48} \\
				\bitbox{1}{\tiny SD} & \bitbox{15}{\tiny WPRI} \\
				\\[-1.5ex] 
				
				\bitbox[]{1}{\tiny 47} & \bitbox[]{12}{} & \bitbox[]{1}{\tiny 34} & \bitbox[]{1}{\tiny 33} & \bitbox[]{1}{\tiny 32} \\
				\bitbox{14}{\tiny WPRI} & \bitbox{2}{\tiny UXL[1:0]} \\
				\\[-1.5ex]
				
				\bitbox[]{1}{\tiny 31} & \bitbox[]{5}{} & \bitbox[]{1}{\tiny 25} & \bitbox[]{1}{\tiny 24} & \bitbox[]{1}{\tiny 23} & \bitbox[]{1}{\tiny 22} & \bitbox[]{1}{} & \bitbox[]{1}{\tiny 20} & \bitbox[]{1}{\tiny 19} & \bitbox[]{1}{\tiny 18} & \bitbox[]{1}{\tiny 17} & \bitbox[]{1}{\tiny 16} \\
				\bitbox{7}{\tiny WPRI} & \bitbox{1}{\tiny SDT} & \bitbox{1}{\tiny SPELP} & \bitbox{3}{\tiny WPRI} & \bitbox{1}{\tiny MXR} & \bitbox{1}{\tiny SUM} & \bitbox{1}{\tiny WPRI} & \bitbox{1}{\tiny XS[1:0]} \\
				\\[-1.5ex]
				
				\bitbox[]{1}{\tiny 15} & \bitbox[]{1}{\tiny 14} & \bitbox[]{1}{\tiny 13} & \bitbox[]{1}{\tiny 12} & \bitbox[]{1}{\tiny 11} & \bitbox[]{1}{\tiny 10} & \bitbox[]{1}{\tiny 9} & \bitbox[]{1}{\tiny 8} & \bitbox[]{1}{\tiny 7} & \bitbox[]{1}{\tiny 6} & \bitbox[]{1}{\tiny 5} & \bitbox[]{1}{\tiny 4} & \bitbox[]{1}{} & \bitbox[]{1}{\tiny 2} & \bitbox[]{1}{\tiny 1} & \bitbox[]{1}{\tiny 0} \\
				\bitbox{1}{\tiny XS[1:0]} & \bitbox{2}{\tiny FS[1:0]} & \bitbox{2}{\tiny WPRI} & \bitbox{2}{\tiny VS[1:0]} & \bitbox{1}{\tiny SPP} & \bitbox{1}{\tiny WPRI} & \bitbox{1}{\tiny UBE} & \bitbox{1}{\tiny SPIE} & \bitbox{3}{\tiny WPRI} & \bitbox{1}{\tiny SIE} & \bitbox{1}{\tiny WPRI}
				
			\end{bytefield}
			\caption{Supervisor-mode status (\texttt{sstatus}) register when $SXLEN=64$.}
			\label{fig:sstatus}
		\end{figure}
	
	\subsubsection{Supervisor Trap Vector Base Address (stvec) Register}
		The stvec register functions as the Supervisor-mode trap vector register, holding the entry point address for S-mode exception handlers.
		Just like its Machine-mode counterpart, it can be configured for either Direct or Vectored trap handling modes.
		This ensures that the OS kernel can catch and process its own traps independently of the higher-level firmware.
		
		\begin{figure}[!htbp]
			\centering
			\begin{bytefield}[bitwidth=0.06\textwidth, bitheight=3ex]{16}
				
				\bitbox[]{1}{\raggedright \tiny SXLEN-1} & \bitbox[]{12}{} & \bitbox[]{1}{\tiny 2} & \bitbox[]{1}{\tiny 1} & \bitbox[]{1}{\tiny 0} \\
				
				\bitbox{14}{\tiny BASE[SXLEN-1:2] (WARL)} & \bitbox{2}{\tiny MODE (WARL)} \\
				
				\bitbox[]{14}{\tiny SXLEN-2} & \bitbox[]{2}{\tiny 2}
				
			\end{bytefield}
			\caption{Supervisor trap vector base address (\texttt{stvec}) register.}
			\label{fig:stvec}
		\end{figure}
	
	\subsubsection{Supervisor Interrupt (sip and sie) Registers}
		The sip register shows which interrupts are currently pending at the Supervisor level, acting as a subset of the Machine-level interrupt status.
		It allows the kernel to check for incoming software or timer interrupts that have been delegated to it by the Machine mode.
		This is the primary way an operating system monitors its own asynchronous events.
		
		\begin{figure}[!htbp]
			\centering
			\begin{bytefield}[bitwidth=0.06\textwidth, bitheight=3ex]{16}
				
				\bitbox[]{2}{\raggedright \tiny SXLEN-1} & \bitbox[]{13}{} & \bitbox[]{1}{\raggedleft \tiny 0} \\
				
				\bitbox{16}{\tiny Interrupts (WARL)} \\
				
				\bitbox[]{16}{\tiny SXLEN}
				
			\end{bytefield}
			\caption{Supervisor interrupt-pending register (\texttt{sip}).}
			\label{fig:sip}
		\end{figure}
		
		The sie register provides individual enable bits for interrupts that are visible to and manageable by the Supervisor level.
		The kernel uses this register to mask or unmask specific interrupt sources according to its current execution priority or state.
		Proper management here is essential for maintaining kernel stability and preventing undesired interrupt recursion.
		
		\begin{figure}[!htbp]
			\centering
			\begin{bytefield}[bitwidth=0.06\textwidth, bitheight=3ex]{16}
				
				\bitbox[]{2}{\raggedright \tiny SXLEN-1} & \bitbox[]{13}{} & \bitbox[]{1}{\raggedleft \tiny 0} \\
				
				\bitbox{16}{\tiny Interrupts (WARL)} \\
				
				\bitbox[]{16}{\tiny SXLEN}
				
			\end{bytefield}
			\caption{Supervisor interrupt-enable register (\texttt{sie}).}
			\label{fig:sie}
		\end{figure}
	
	\subsubsection{Counter-Enable (scounteren) Register}
		The scounteren register allows Supervisor mode to control whether User-mode applications can access hardware performance counters.
		It mirrors the functionality of the Machine-mode version, letting the OS selectively expose metrics like instruction counts to user programs.
		This provides a layered security approach to performance monitoring within the application environment.
		
		\begin{figure}[!htbp]
			\centering
			\begin{bytefield}[bitwidth=0.06\textwidth, bitheight=3ex]{16}
				
				\bitbox[]{1}{\tiny 31} & \bitbox[]{1}{\tiny 30} & \bitbox[]{1}{\tiny 29} & \bitbox[]{1}{\tiny 28} & \bitbox[]{5}{} & \bitbox[]{1}{\tiny 6} & \bitbox[]{1}{\tiny 5} & \bitbox[]{1}{\tiny 4} & \bitbox[]{1}{\tiny 3} & \bitbox[]{1}{\tiny 2} & \bitbox[]{1}{\tiny 1} & \bitbox[]{1}{\tiny 0} \\
				
				\bitbox{1}{\tiny HPM31} & \bitbox{1}{\tiny HPM30} & \bitbox{1}{\tiny HPM29} & \bitbox{7}{\tiny ...} & \bitbox{1}{\tiny HPM5} & \bitbox{1}{\tiny HPM4} & \bitbox{1}{\tiny HPM3} & \bitbox{1}{\tiny IR} & \bitbox{1}{\tiny TM} & \bitbox{1}{\tiny CY} \\
				
				\bitbox[]{1}{\tiny 1} & \bitbox[]{1}{\tiny 1} & \bitbox[]{1}{\tiny 1} & \bitbox[]{7}{\tiny 23} & \bitbox[]{1}{\tiny 1} & \bitbox[]{1}{\tiny 1} & \bitbox[]{1}{\tiny 1} & \bitbox[]{1}{\tiny 1} & \bitbox[]{1}{\tiny 1} & \bitbox[]{1}{\tiny 1}
				
			\end{bytefield}
			\caption{Counter-enable (\texttt{scounteren}) register.}
			\label{fig:scounteren}
		\end{figure}
	
	\subsubsection{Supervisor Scratch (sscratch) Register}
		The sscratch register is a dedicated scratch register used by the Supervisor-mode trap handler to facilitate quick state saving.
		The kernel typically stores a pointer to per-CPU data or a kernel stack here to avoid using general-purpose registers before they are saved.
		This provides a vital "bridge" to a safe execution environment when a trap occurs in user space.
		
		\begin{figure}[!htbp]
			\centering
			\begin{bytefield}[bitwidth=0.06\textwidth, bitheight=3ex]{16}
				
				\bitbox[]{2}{\raggedright \tiny SXLEN-1} & \bitbox[]{13}{} & \bitbox[]{1}{\raggedleft \tiny 0} \\
				
				\bitbox{16}{\tiny sscratch} \\
				
				\bitbox[]{16}{\tiny SXLEN}
				
			\end{bytefield}
			\caption{Supervisor Scratch Register.}
			\label{fig:sscratch}
		\end{figure}
	
	\subsubsection{Supervisor Exception Program Counter (sepc) Register}
		The sepc register automatically saves the address of the instruction that was interrupted when a trap enters Supervisor mode.
		It provides the return address needed by the sret instruction to resume normal program execution after the kernel finishes its task.
		It is the S-mode functional equivalent to the Machine-mode exception program counter.
		
		\begin{figure}[!htbp]
			\centering
			\begin{bytefield}[bitwidth=0.06\textwidth, bitheight=3ex]{16}
				
				\bitbox[]{2}{\raggedright \tiny SXLEN-1} & \bitbox[]{13}{} & \bitbox[]{1}{\raggedleft \tiny 0} \\
				
				\bitbox{16}{\tiny sepc} \\
				
				\bitbox[]{16}{\tiny SXLEN}
				
			\end{bytefield}
			\caption{Supervisor exception program counter register.}
			\label{fig:sepc}
		\end{figure}
	
	\subsubsection{Supervisor Cause (scause) Register}
		The scause register stores the interrupt or exception code for traps that are handled within Supervisor mode.
		It is the central piece of data used by the kernel's trap entry point to dispatch control to the specific handler function.
		It allows the OS to distinguish between system calls from user space and various hardware interrupts.
		
		\begin{figure}[!htbp]
			\centering
			\begin{bytefield}[bitwidth=0.06\textwidth, bitheight=3ex]{16}
				
				\bitbox[]{2}{\raggedright \tiny SXLEN-1} & \bitbox[]{2}{\raggedright \tiny SXLEN-2} & \bitbox[]{11}{} & \bitbox[]{1}{\raggedleft \tiny 0} \\
				
				\bitbox{2}{\tiny Interrupt} & \bitbox{14}{\tiny Exception Code (WLRL)} \\
				
				\bitbox[]{2}{\tiny 1} & \bitbox[]{14}{\tiny SXLEN-1}
				
			\end{bytefield}
			\caption{Supervisor Cause (\texttt{scause}) register.}
			\label{fig:scause}
		\end{figure}
	
	\subsubsection{Supervisor Trap Value (stval) Register}
		The stval register holds auxiliary data for Supervisor-mode traps, such as the specific virtual address that caused a page fault.
		This information is indispensable for the kernel's memory management system to resolve faults or log access errors.
		It provides the necessary context to handle user-space exceptions effectively without requiring Machine-mode intervention.
		
		\begin{figure}[!htbp]
			\centering
			\begin{bytefield}[bitwidth=0.06\textwidth, bitheight=4ex]{16}
				
				\bitbox[]{2}{\raggedright \tiny SXLEN-1} & \bitbox[]{13}{} & \bitbox[]{1}{\raggedleft \tiny 0} \\
				
				\bitbox{16}{\tiny stval} \\
				
				\bitbox[]{16}{\tiny SXLEN}
				
			\end{bytefield}
			\caption{Supervisor Trap Value register.}
			\label{fig:stval}
		\end{figure}